%%% =======================================================================
%%% ISLS Extension Phase 12 v1.0.0
%%% The Agent: Studio as Workspace-Aware Development Exoskeleton
%%% The Lab That Makes the Product Pure
%%% Author: Sebastian Klemm
%%% Date: 19 March 2026
%%% Status: Implementation Specification for Claude Code
%%% =======================================================================
\documentclass[11pt,a4paper,twoside]{article}

\usepackage[utf8]{inputenc}
\usepackage{amsmath,amssymb,amsthm}
\usepackage{booktabs,longtable,tabularx}
\usepackage{enumitem}
\usepackage{geometry}
\usepackage{hyperref}
\usepackage{xcolor}
\usepackage{fancyhdr}
\usepackage{listings}
\usepackage{titlesec}

\geometry{margin=2.5cm,top=3cm,bottom=3cm,headheight=14pt}

\theoremstyle{definition}
\newtheorem{definition}{Definition}[section]
\newtheorem{requirement}{Requirement}[section]

\theoremstyle{remark}
\newtheorem{remark}{Remark}[section]

\newcommand{\ISLS}{\textsf{ISLS}}

\definecolor{sec-color}{HTML}{1e3a5f}

\titleformat{\section}
  {\Large\bfseries\color{sec-color}}
  {\thesection}{1em}{}

\lstset{
  basicstyle=\ttfamily\small,
  breaklines=true,
  frame=single,
  backgroundcolor=\color{gray!5},
  rulecolor=\color{gray!30},
}

\pagestyle{fancy}
\fancyhf{}
\fancyhead[LE,RO]{\ISLS{} Phase 12 v1.0.0}
\fancyhead[RE,LO]{Sebastian Klemm}
\fancyfoot[C]{\thepage}

\hypersetup{
  colorlinks=true,
  linkcolor=blue!60!black,
  pdftitle={ISLS Phase 12 — The Agent},
  pdfauthor={Sebastian Klemm},
}

\begin{document}

\begin{titlepage}
\centering
\vspace*{2.5cm}
{\Huge\bfseries \ISLS{} Extension Phase~12\\[0.3cm]
v1.0.0}\\[1.5cm]
{\Large The Agent}\\[1cm]
{\large The Studio becomes the workshop.\\
Chat-driven. Workspace-aware. Memory-backed.\\
Formally verified. Accumulating. Autonomous.\\[0.5cm]
Everyone has the same LLM.\\
Nobody has this lab.}\\[2cm]
{\large Sebastian Klemm}\\[0.5cm]
{\large 19 March 2026}\\[1.5cm]

\begin{abstract}
\noindent
This document specifies the transformation of the \ISLS{} Studio
from a dashboard into an \textbf{agent}: a workspace-aware,
chat-driven development interface that uses any LLM as its engine
but wraps it in the full \ISLS{} substrate --- Pattern Memory,
Constraint Propagation, Foundry compile-test-fix, formal
crystallization, and accumulating autonomy.

The operator opens \texttt{localhost:8420}, types what they want in
natural language, and the Agent does the rest: reads the project,
understands what exists, constructs a maximally-constrained prompt,
calls the LLM, receives code, writes it to disk, compiles, tests,
fixes if needed, crystallizes, stores the pattern. Next time
something similar is needed, it comes from memory.

This is not a Claude Code replacement. Claude Code is the LLM.
\ISLS{} is the exoskeleton. The exoskeleton makes the LLM's
output reproducible, verifiable, accumulating, and progressively
cheaper. Everyone has access to the same LLM API. Nobody else has
this infrastructure around it.

The value proposition in one sentence: \textbf{the 100th project
costs less than the 1st, and the 1000th costs almost nothing.}

\medskip\noindent\textbf{Status:} Implementation Specification.\\
\textbf{Modifies:} C19 (Gateway/Studio), C25 (Oracle), C27 (Foundry).\\
\textbf{New module:} \texttt{agent.rs} in \texttt{isls-gateway}.\\
\textbf{No new crate.} The Agent is a mode of the Studio, not a
separate system.
\end{abstract}
\end{titlepage}

\tableofcontents
\newpage


%%% =====================================================================
\part{The Problem}\label{part:problem}

\section{Why ISLS Is Currently Useless for Daily Work}\label{sec:useless}

Today, to use the Forge, the operator must:
\begin{enumerate}[nosep]
  \item Write a JSON DecisionSpec by hand.
  \item Run \texttt{isls forge --spec file.json --domain rust}.
  \item Wait for output.
  \item Manually inspect the output.
  \item Manually copy files into their project.
  \item Manually run \texttt{cargo check} to see if it works.
  \item If it doesn't work, start over.
\end{enumerate}

Claude Code, by comparison: the operator types what they want, it
reads the project, edits files, compiles, tests, iterates, done.
No JSON specs, no manual copying, no separate tools.

The gap is not intelligence. The gap is \textbf{agency} ---
the ability to act on the operator's behalf in the operator's
project, autonomously, in a feedback loop.


%%% =====================================================================
\part{The Agent}\label{part:agent}

\section{What It Looks Like}\label{sec:looks}

\begin{lstlisting}[title={The Studio Agent Interface}]
+----------------------------------------------------------+
| ISLS Agent                    [project: ~/my-api]  [M:23]|
+----------------------------------------------------------+
|                                                          |
|  YOU: Add a /search endpoint that does full-text         |
|       search on bookmarks by title and tags              |
|                                                          |
|  AGENT: Analyzing workspace...                           |
|    Found: src/router.rs (6 routes), src/models.rs        |
|    (Bookmark struct with title, url, tags fields),       |
|    src/database.rs (SQLx, bookmarks table)               |
|                                                          |
|    Plan:                                                 |
|    1. Add SearchParams struct to models.rs               |
|    2. Add search_bookmarks() to database.rs              |
|    3. Add GET /search route to router.rs                 |
|    4. Add test in tests/search.rs                        |
|                                                          |
|    Executing...                                          |
|    [1/4] models.rs: SearchParams added (from memory)     |
|    [2/4] database.rs: search query (oracle: 0.8s)        |
|    [3/4] router.rs: route added (from memory)            |
|    [4/4] tests/search.rs: test generated (oracle: 0.6s)  |
|                                                          |
|    > cargo check ... OK                                  |
|    > cargo test  ... 13/13 passed                        |
|                                                          |
|    Crystal: a7b3c9 | Files: 4 | Attempts: 1             |
|    Patterns stored: 3 new, 1 from memory                 |
|                                                          |
|  YOU: The search should also support date ranges          |
|                                                          |
|  AGENT: Analyzing diff...                                |
|    Modifying: SearchParams (add from/to dates),          |
|    search_bookmarks() (add WHERE clause),                |
|    test (add date range test case)                       |
|    ...                                                   |
+----------------------------------------------------------+
| [M] Memory: 23 patterns | Auto: 34% | Cost: $0.02       |
+----------------------------------------------------------+
\end{lstlisting}


\section{The Five Capabilities}\label{sec:capabilities}

\subsection{1. Chat Interface}\label{sec:chat}

\begin{definition}[Agent Chat]\label{def:chat}
The Studio gains a primary chat view. The operator types natural
language. The Agent responds with actions. The conversation is
persistent within a session and can reference previous messages.

\begin{lstlisting}
// WebSocket message types (extend existing ws.rs)
pub enum AgentMessage {
    /// Operator sends a request
    UserMessage { content: String },

    /// Agent shows analysis
    Analysis { workspace_summary: String, plan: Vec<PlanStep> },

    /// Agent shows progress
    Progress { step: usize, total: usize, description: String,
               source: String }, // "memory", "oracle", "static"

    /// Agent shows compile/test result
    ToolResult { tool: String, success: bool, output: String },

    /// Agent shows final result
    Complete { files_changed: Vec<String>, crystal_id: String,
               patterns_stored: usize, patterns_reused: usize,
               cost_usd: f64 },

    /// Agent asks for clarification
    Clarification { question: String, options: Vec<String> },
}
\end{lstlisting}
\end{definition}


\subsection{2. Workspace Awareness}\label{sec:workspace}

\begin{definition}[Project Binding]\label{def:projectbind}
The Agent is always bound to a project directory. At startup or
via the UI, the operator selects a directory. The Agent runs
WorkspaceAnalyzer (C27) and maintains a live model:

\begin{lstlisting}
pub struct AgentWorkspace {
    pub root: PathBuf,
    pub model: WorkspaceModel,      // from C27 WorkspaceAnalyzer
    pub file_index: FileIndex,       // path -> last_modified + summary
    pub type_index: TypeIndex,       // type_name -> file + fields
    pub function_index: FunctionIndex, // fn_name -> file + signature
    pub route_index: RouteIndex,     // path -> handler + method (if API)
    pub last_analysis: Instant,
}

impl AgentWorkspace {
    /// Re-analyze when files change
    pub fn refresh_if_stale(&mut self) -> bool;

    /// Find all files relevant to a task
    pub fn relevant_files(&self, task: &str) -> Vec<RelevantFile>;

    /// Get the full content of relevant files for the Oracle prompt
    pub fn context_for_prompt(&self, relevant: &[RelevantFile],
                               max_tokens: usize) -> String;

    /// After the Agent modifies files, update indexes
    pub fn notify_change(&mut self, paths: &[PathBuf]);
}
\end{lstlisting}
\end{definition}

\begin{definition}[Context Construction]\label{def:context}
The Agent builds the Oracle prompt from workspace context:
\begin{lstlisting}
pub fn build_agent_prompt(
    user_message: &str,
    workspace: &AgentWorkspace,
    memory_matches: &[PatternEntry],
    conversation_history: &[ConversationTurn],
) -> SynthesisPrompt {
    let relevant = workspace.relevant_files(user_message);
    let context = workspace.context_for_prompt(&relevant, 6000);

    let system = format!(
        "You are a code agent working in a Rust project.\n\
         You modify EXISTING files. You do NOT create from scratch.\n\
         Output ONLY the complete modified file content.\n\
         No explanations, no markdown fences.\n\n\
         PROJECT STRUCTURE:\n{}\n\n\
         EXISTING CODE (relevant files):\n{}\n\n\
         KNOWN PATTERNS:\n{}",
        workspace.model.summary(),
        context,
        format_memory_matches(memory_matches),
    );

    let user = format!(
        "CONVERSATION:\n{}\n\nLATEST REQUEST:\n{}",
        format_history(conversation_history),
        user_message,
    );

    SynthesisPrompt {
        system,
        user,
        output_format: OutputFormat::Rust,
        max_tokens: 4096,
        temperature: 0.0,
    }
}
\end{lstlisting}

The key difference from the current Forge: the prompt contains the
\textbf{actual code} from the operator's project. The LLM sees what
exists and modifies it, rather than generating from scratch.
\end{definition}


\subsection{3. Autonomous Execution}\label{sec:execution}

\begin{definition}[Agent Execution Loop]\label{def:execloop}
\begin{lstlisting}
pub struct Agent {
    workspace: AgentWorkspace,
    oracle: Box<dyn SynthesisOracle>,
    foundry: Foundry,
    memory: PatternMemory,
    conversation: Vec<ConversationTurn>,
    config: AgentConfig,
}

impl Agent {
    pub async fn handle_message(&mut self, msg: &str) -> AgentResult {
        // 1. Analyze: what does the operator want?
        let task = self.analyze_intent(msg).await?;

        // 2. Plan: which files need to change?
        let plan = self.plan_changes(&task).await?;
        self.emit(AgentMessage::Analysis { ... });

        // 3. For each planned change:
        let mut changed_files = Vec::new();
        for (i, step) in plan.steps.iter().enumerate() {
            // 3a. Check Pattern Memory first
            let signature = step.to_five_d_signature();
            if let Some(pattern) = self.memory.find_match(&signature, 0.85) {
                // Adapt pattern to current context
                let adapted = self.adapt_pattern(&pattern, &step)?;
                self.write_file(&step.file, &adapted)?;
                self.emit(AgentMessage::Progress {
                    source: "memory".into(), ...
                });
            } else {
                // 3b. Call Oracle with full workspace context
                let prompt = build_agent_prompt(
                    &step.description,
                    &self.workspace,
                    &self.memory.recent_matches(5),
                    &self.conversation,
                );
                let response = self.oracle.synthesize(&prompt)?;

                // 3c. Strip markdown, extract code
                let code = strip_and_extract(&response.content);
                self.write_file(&step.file, &code)?;

                self.emit(AgentMessage::Progress {
                    source: "oracle".into(), ...
                });
            }
            changed_files.push(step.file.clone());
        }

        // 4. Compile and test
        let compile = self.foundry.toolchain.cargo_check(&self.workspace.root)?;
        self.emit(AgentMessage::ToolResult {
            tool: "cargo check".into(), ...
        });

        if !compile.success {
            // 5. Auto-fix: send error back to Oracle
            let fix_result = self.auto_fix(&compile.stderr, &changed_files).await?;
            if !fix_result.success {
                return Ok(AgentResult::CompileFailure { ... });
            }
        }

        let test = self.foundry.toolchain.cargo_test(&self.workspace.root)?;
        self.emit(AgentMessage::ToolResult {
            tool: "cargo test".into(), ...
        });

        // 6. Crystallize
        let crystal = self.crystallize(&changed_files, &task)?;

        // 7. Store patterns for reuse
        let patterns_stored = self.store_patterns(&changed_files, &crystal)?;

        // 8. Update workspace model
        self.workspace.notify_change(&changed_files);

        // 9. Update conversation
        self.conversation.push(ConversationTurn {
            role: "user".into(), content: msg.into(),
        });
        self.conversation.push(ConversationTurn {
            role: "agent".into(), content: format!("Changed {} files", changed_files.len()),
        });

        Ok(AgentResult::Success { crystal, patterns_stored, ... })
    }
}
\end{lstlisting}
\end{definition}


\subsection{4. Multi-Turn Conversation}\label{sec:multiturn}

\begin{definition}[Conversation Context]\label{def:conversation}
The Agent maintains conversation history within a session. Each
turn is stored and included in subsequent Oracle prompts. This
enables:

\begin{itemize}[nosep]
  \item \textbf{Refinement}: ``The search should also support date
        ranges'' --- the Agent sees the previous turn (search endpoint)
        and modifies only what changed.
  \item \textbf{Reference}: ``Make the same change in the other module''
        --- the Agent understands ``the same change'' from context.
  \item \textbf{Correction}: ``That's wrong, the field is called
        \texttt{created\_at} not \texttt{date}'' --- the Agent re-reads
        the file and corrects.
\end{itemize}

Conversation history is truncated to fit the LLM context window
(keep last $n$ turns, summarize earlier ones).
\end{definition}


\subsection{5. Accumulating Advantage}\label{sec:accumulate}

\begin{definition}[The Compounding Effect]\label{def:compound}
Every Agent action that succeeds is crystallized and stored in
Pattern Memory. The accumulation metrics:

\begin{lstlisting}
pub struct AccumulationMetrics {
    pub total_requests: u64,
    pub memory_served: u64,          // no LLM call needed
    pub oracle_served: u64,          // LLM called
    pub oracle_first_try: u64,      // compiled on first attempt
    pub oracle_with_fix: u64,        // needed correction cycle
    pub autonomy_ratio: f64,         // memory_served / total
    pub total_cost_usd: f64,
    pub avg_cost_per_request: f64,   // should decrease over time
    pub patterns_in_memory: usize,
}
\end{lstlisting}

The Studio dashboard shows:
\begin{itemize}[nosep]
  \item Autonomy ratio as a percentage bar (``34\% self-sufficient'')
  \item Cost trend chart (cost per request declining over time)
  \item Pattern count growing
  \item ``Saved \$X.XX by using memory instead of LLM''
\end{itemize}
\end{definition}


%%% =====================================================================
\part{Intent Analysis}\label{part:intent}

\section{What the Agent Understands}\label{sec:understands}

\begin{definition}[Intent Classification]\label{def:intent}
The Agent classifies the operator's message into an operation type.
This classification can be done with simple keyword matching first
(no LLM needed), or by a lightweight LLM call for ambiguous cases:

\begin{lstlisting}
pub enum AgentIntent {
    /// "Add a search endpoint" -> add new code
    Add { target: String, description: String },

    /// "Fix the compile error in database.rs" -> fix existing code
    Fix { file: Option<String>, issue: String },

    /// "Refactor the handler into smaller functions" -> restructure
    Refactor { target: String, instruction: String },

    /// "Write tests for the search module" -> generate tests
    Test { target: Option<String> },

    /// "What does the search function do?" -> explain (no code change)
    Explain { target: String },

    /// "Run the tests" -> execute tool
    Run { command: String },

    /// "Show me the project structure" -> workspace query
    Query { question: String },

    /// Complex: needs LLM to decompose into sub-intents
    Complex { description: String },
}
\end{lstlisting}

For \texttt{Add}, \texttt{Fix}, \texttt{Refactor}, \texttt{Test}:
the Agent acts (modifies files, compiles, tests).

For \texttt{Explain}, \texttt{Query}: the Agent responds without
changing files.

For \texttt{Run}: the Agent executes a shell command and shows
output.

For \texttt{Complex}: the Agent decomposes into a multi-step plan
(each step is one of the simple intents above).
\end{definition}


%%% =====================================================================
\part{File Operations}\label{part:fileops}

\section{How the Agent Modifies Files}\label{sec:modify}

\begin{definition}[Targeted File Modification]\label{def:modify}
The Agent does NOT generate entire files from scratch (wasteful,
error-prone). It generates \textbf{patches}: targeted modifications
to existing files.

Two strategies:
\begin{enumerate}[nosep]
  \item \textbf{Full-file replacement}: for small files or new files,
        the Oracle returns the complete file content. The Agent writes
        it directly.
  \item \textbf{Surgical edit}: for large files, the Agent sends only
        the relevant section to the Oracle, receives the modified
        section, and splices it back in. This reduces token usage and
        error surface.
\end{enumerate}

\begin{lstlisting}
pub enum FileOperation {
    Create { path: String, content: String },
    Replace { path: String, content: String },
    Edit {
        path: String,
        anchor: String,      // unique string to locate the edit point
        old_content: String, // what to replace
        new_content: String, // replacement
    },
    Append { path: String, content: String },
    Delete { path: String },
}
\end{lstlisting}

For surgical edits, the Oracle prompt includes:
\begin{lstlisting}
"File: src/database.rs
REPLACE THIS:
    pub async fn list_bookmarks(...) -> ... {
        sqlx::query_as!(...)
    }
WITH YOUR MODIFICATION.
Context: adding full-text search on title and tags fields.
Output ONLY the replacement code block. Nothing else."
\end{lstlisting}
\end{definition}


%%% =====================================================================
\part{Gateway Integration}\label{part:gateway}

\section{New Endpoints}\label{sec:endpoints}

\begin{lstlisting}
// Agent-specific endpoints
POST  /agent/message        Body: {content: "add search endpoint"}
                            Returns: stream of AgentMessage via SSE

POST  /agent/bind           Body: {project_dir: "/path/to/project"}
                            Returns: WorkspaceModel summary

GET   /agent/workspace      Returns: current workspace model

GET   /agent/conversation   Returns: conversation history

POST  /agent/clear          Clears conversation history

GET   /agent/accumulation   Returns: AccumulationMetrics

// The chat interface uses Server-Sent Events (SSE) for streaming:
// POST /agent/message returns a stream of AgentMessage as they occur
// (analysis, progress, tool results, completion)
\end{lstlisting}


\section{Studio Chat View}\label{sec:chatview}

The Studio's primary view becomes the Chat. The 7 existing views
remain accessible via sidebar, but the Chat is default.

\begin{lstlisting}[title={Studio Layout with Chat}]
+--[sidebar]--+--[chat]------------------------------+
| [C] Chat    | YOU: add search endpoint              |
| [D] Dash    |                                       |
| [F] Forge   | AGENT: Analyzing...                   |
| [E] Explore |   Found: router.rs, models.rs, ...    |
| [M] Monitor |   Plan: 4 files                       |
| [B] Build   |   [1/4] models.rs (memory) OK         |
| [O] Oracle  |   [2/4] database.rs (oracle 0.8s) OK  |
| [C] Const.  |   [3/4] router.rs (memory) OK         |
| [N] Nav     |   [4/4] tests/search.rs (oracle) OK   |
|             |   > cargo check ... OK                |
|             |   > cargo test ... 13/13 passed       |
|             |   Crystal: a7b3c9                     |
|             |                                       |
| Auto: 34%   | YOU: also support date ranges         |
| Mem: 23     |                                       |
| $0.02       | AGENT: Modifying 3 files...            |
+-----------  +---------------------------------------+
\end{lstlisting}


%%% =====================================================================
\part{Configuration}\label{part:config}

\begin{lstlisting}
agent:
  enabled: true
  default_project: "~/projects/my-api"
  max_conversation_turns: 20
  max_context_tokens: 6000     # workspace context in prompt
  auto_compile: true           # compile after every change
  auto_test: true              # test after every change
  auto_crystallize: true       # crystallize successful changes
  auto_store_patterns: true    # store patterns from successful changes
  surgical_edit_threshold: 200 # lines; above this use surgical edit
  oracle_provider: "auto"      # auto-detect from env vars
\end{lstlisting}


%%% =====================================================================
\part{Acceptance Tests}\label{part:tests}

\begin{longtable}{l l p{5.5cm}}
\toprule
\textbf{ID} & \textbf{Name} & \textbf{Procedure} \\
\midrule
\endfirsthead \midrule \endhead \bottomrule \endfoot
AT-AG1 & Project bind & Bind to a Rust project; verify workspace
  model contains correct modules, types, functions. \\
AT-AG2 & Intent classification & Send ``add an endpoint''; verify
  classified as Add. Send ``fix the bug''; verify Fix. \\
AT-AG3 & Workspace context & Bind project with 5 files; verify
  prompt contains relevant file contents. \\
AT-AG4 & File creation & Send ``create a new module health.rs'';
  verify file created on disk. \\
AT-AG5 & File modification & Send ``add a field to User struct'';
  verify existing file modified correctly. \\
AT-AG6 & Compile check & Agent modifies file; verify
  \texttt{cargo check} automatically runs. \\
AT-AG7 & Auto-fix & Mock Oracle returns code with missing import;
  verify Agent detects compile error and retries with error context. \\
AT-AG8 & Crystallization & Successful change produces a crystal
  with evidence chain linking to the conversation. \\
AT-AG9 & Pattern storage & After successful change, verify pattern
  stored in memory with correct signature. \\
AT-AG10 & Memory reuse & Store a pattern; request similar change;
  verify served from memory without Oracle call. \\
AT-AG11 & Multi-turn & Send ``add endpoint'', then ``also add
  pagination''; verify second turn modifies the first result. \\
AT-AG12 & Accumulation metrics & Make 5 requests (3 oracle, 2 memory);
  verify autonomy ratio = 40\%. \\
AT-AG13 & SSE streaming & Connect to /agent/message; verify
  AgentMessages stream in real-time. \\
AT-AG14 & Explain intent & Send ``what does search\_bookmarks do?'';
  verify response without file changes. \\
\end{longtable}


%%% =====================================================================
\part{Handoff}\label{part:handoff}

\section{What Changes}

\begin{tabularx}{\textwidth}{l X}
\toprule
\textbf{File} & \textbf{Change} \\
\midrule
\texttt{isls-gateway/src/agent.rs} & New: Agent struct, execution loop,
  intent classification, workspace binding. \\
\texttt{isls-gateway/src/lib.rs} & Add agent routes: /agent/message,
  /agent/bind, /agent/workspace, /agent/conversation,
  /agent/accumulation. \\
\texttt{isls-gateway/src/ws.rs} & Add AgentMessage event types for
  SSE streaming. \\
\texttt{isls-gateway/src/static/studio.html} & Add Chat view as
  primary view. SSE client for streaming agent responses. \\
\texttt{isls-oracle/src/lib.rs} & Expose \texttt{strip\_markdown\_fences}
  as public utility. \\
\texttt{isls-foundry/src/workspace.rs} & Expose WorkspaceAnalyzer
  fields for Agent consumption (type\_index, function\_index). \\
\texttt{isls-cli/src/main.rs} & Add \texttt{isls agent} subcommand
  for terminal-based agent (non-GUI fallback). \\
\bottomrule
\end{tabularx}

\section{Test Count}

\begin{tabular}{l r r}
\toprule
\textbf{Phase} & \textbf{New Tests} & \textbf{Cumulative} \\
\midrule
Phase 1--11 & 373 & 373 \\
Phase 12 (Agent) & 14 & $\ge$ 387 \\
\bottomrule
\end{tabular}

\section{Completion Criteria}

\begin{enumerate}[nosep]
  \item Operator types natural language in Studio Chat.
  \item Agent analyzes workspace and shows plan before executing.
  \item Agent modifies existing files with workspace context.
  \item \texttt{cargo check} and \texttt{cargo test} run automatically.
  \item Compile errors trigger auto-fix cycle (error $\to$ Oracle $\to$
        retry).
  \item Every successful change is crystallized with evidence chain.
  \item Patterns are stored and reused for similar future requests.
  \item Accumulation metrics show autonomy ratio, cost trend, pattern
        count.
  \item Multi-turn conversation works (refinement, correction,
        reference).
  \item SSE streaming shows real-time progress in the Studio.
  \item Terminal fallback: \texttt{isls agent} works without browser.
  \item Total tests $\ge 387$, zero warnings.
  \item \textbf{Everyone has the same LLM. Nobody has this lab.}
\end{enumerate}


\vfill
\begin{center}
\rule{0.5\textwidth}{0.4pt}\\[0.5cm]
{\small Generated for \ISLS{} Extension Phase 12 v1.0.0.\\[0.2cm]
The 1st project costs \$0.05 and takes 10 seconds.\\
The 100th project costs \$0.01 and takes 3 seconds.\\
The 1000th project costs nothing and takes milliseconds.\\[0.3cm]
The LLM is the chemist.\\
The Memory is the recipe book that grows.\\
The Foundry is the quality control.\\
The Crystal is the proof of purity.\\
The Lab is the competitive advantage.\\[0.3cm]
The chemist is for hire. The lab is yours.}
\end{center}

\end{document}
